% Recuperación del producto nativo y su selector, no incorporación de RH.
\subsection{Irréductibilité multiplicative et périodes de retour}
\label{primos-retornos}

La distinction entre structure, évaluation et phase apparaît également dans le
secteur multiplicatif. Le produit local d’APP contient tant le résidu
que le quotient ; son itération permet de construire des formes normales et des
factorisations avant d’examiner leurs périodes décimales. Ce développement
situe l’expression \emph{mode premier} dans un domaine précis et relie
l’irréductibilité aux opérations déjà définies, non à une coïncidence
isolée de périodes (voir \ref{primos-retornos}).

Soit
\[
 \mathcal W_9=\{\varnothing\}\sqcup
 \bigsqcup_{\ell\geq1}\{1,\ldots,9\}^{\ell}.
\]
Les suites sont ordonnées à partir du chiffre le moins significatif. L'évaluation
\(\operatorname{ev}_9(u)=\sum_j u_j9^j\), avec
\(\operatorname{ev}_9(\varnothing)=0\), est bijective sur les entiers non
négatifs : pour un entier positif, la division de \(n-1\) par \(9\)
détermine de manière unique le premier chiffre et la continuation. C'est la
numération nonadique bijective ; elle conserve le représentant \(9\) de la
frontière.

La somme \(\boxplus_\#\) s'obtient en appliquant le feuillet additif d'APP aux chiffres de chaque position et en propageant leurs quotients. S'il n'y a qu'un chiffre, celui-ci constitue le résidu provisoire ; un quotient entrant non nul est normalisé au moyen d'une seconde opération additive. La somme des quotients est transportée à la position suivante. Pour multiplier une suite par un chiffre \(d\), la cellule multiplicative produit \(u_jd=9q_j+r_j\), et le quotient entrant \(c_j\) est incorporé au moyen de
\[
 (w_j,c_{j+1})=
 \begin{cases}
 (r_j,q_j),&c_j=0,\\
 \bigl(R^+(r_j,c_j),q_j+q^+(r_j,c_j)\bigr),&c_j>0.
 \end{cases}
\]
On commence avec \(c_0=0\) et on incorpore le quotient final lorsqu'il est non
nul. Pendant l'addition, \(c_j\leq2\) ; pendant la multiplication par un chiffre,
\(c_j\leq9\). Les opérations locales reçoivent ainsi des arguments dans
leur domaine. Notons \(d\odot_\#u\) cette seconde procédure,
avec \(d\odot_\#\varnothing=\varnothing\).

Pour \(v=(v_0,\ldots,v_{m-1})\), le produit complet est construit par
la récurrence
\[
 a_m=\varnothing,\qquad
 a_j=(9\odot_\#a_{j+1})\boxplus_\#(v_j\odot_\#u),
 \qquad u\boxtimes_\#v=a_0.
\]
La règle compose les deux feuillets locaux. Son évaluation vérifie
\begin{equation}
 \operatorname{ev}_9(u\boxtimes_\#v)
   =\operatorname{ev}_9(u)\operatorname{ev}_9(v).
 \label{eq:primos-entrelazamiento}
\end{equation}
En effet, chaque transition conserve la somme partielle augmentée du quotient
en attente pondéré par la puissance de \(9\). La récurrence donne alors
\(\operatorname{ev}_9(a_j)=9\operatorname{ev}_9(a_{j+1})+
v_j\operatorname{ev}_9(u)\), dont l'itération prouve
\eqref{eq:primos-entrelazamiento}. L'unicité de la forme normale transporte
l'associativité et la commutativité vers le produit construit ; son unité est \((1)\).

\begin{definition}[Coproduit multiplicatif réduit]
Sur l'espace vectoriel à support fini engendré par les formes normales non vides, on définit
\[
 \widetilde\Delta_\# e_w
 =\sum_{\substack{u\boxtimes_\#v=w\\u,v\ne(1)}}e_u\otimes e_v.
\]
Une forme normale non unitaire est irréductible lorsque son image par
\(\widetilde\Delta_\#\) est nulle.
\end{definition}

Les fibres de factorisation sont finies d'après
\eqref{eq:primos-entrelazamiento}. Dans la complétion
\(\mathcal H_\#=\ell^2(\mathcal W_9\setminus\{\varnothing\})\),
les images de formes distinctes ont des supports disjoints. Si
\(d_\#(w)\) compte les factorisations ordonnées non unitaires, le domaine
de la clôture est
\[
 \operatorname{Dom}(\overline{\widetilde\Delta_\#})
 =\left\{x\in\mathcal H_\#:
       \sum_w d_\#(w)|x_w|^2<\infty\right\}.
\]
En effet, le carré de la norme de l'image est cette somme pondérée. Par conséquent, l'opérateur
\[
 A_\#=(\overline{\widetilde\Delta_\#})^*
                  \overline{\widetilde\Delta_\#}
\]
est diagonal, positif et autoadjoint : sa valeur propre en \(e_w\) est le
nombre de factorisations ordonnées non unitaires de \(w\). La projection
\[
 P_{\rm irr}=(I-|e_{(1)}\rangle\langle e_{(1)}|)
                    \mathbf1_{\{0\}}(A_\#)
\]
sélectionne exactement les formes irréductibles. Par l'évaluation multiplicative, celles-ci correspondent aux nombres premiers ordinaires. La sélection dépend de la structure de factorisation et maintient séparé l'état unitaire.

Pour une valeur \(n\geq2\) première avec \(10\), la division décimale induit la
permutation des résidus \(r\mapsto10r\pmod n\). Le retour du résidu
\(1\) a pour période
\[
 \tau_n=\operatorname{ord}_n(10).
\]
Si \(n=\prod_p p^{a_p}\), le théorème chinois des restes donne
\[
 \tau_n=\operatorname{mcm}_{p^{a_p}\parallel n}
                    \operatorname{ord}_{p^{a_p}}(10).
\]
Le retour composé synchronise les retours des puissances premières.
Pour un nombre premier \(p\notin\{2,3,5\}\), on a
\[
 \tau_p\mid p-1,\qquad
 M_p=\frac{10^{\tau_p}-1}{p}\in\mathbb N,
 \qquad M_p\equiv0\pmod9.
\]
La dernière congruence exprime la clôture nonadique de la répétition décimale : \(10^{\tau_p}-1\) est un multiple de \(9\) et \(p\) est inversible modulo \(9\). La même congruence vaut pour tout \(n\) premier avec \(90\). En particulier, \(7,13,77,91,143\) ont pour période décimale \(6\), bien que seuls les deux premiers soient irréductibles. Le coproduit et la période enregistrent donc des propriétés différentes d'une même valeur construite.

Cette différenciation précise la relation avec le raffinement apériodique.
La suite des vacances détermine les incréments de l'indice de
capacité ; la factorisation détermine les modes multiplicatifs ; la
division décimale détermine leurs temps de retour. Le représentant
nonadique conserve la fermeture de phase, tandis que le quotient et la
factorisation conservent une information que cette fermeture ne distingue pas. Le
lien avec les constructions spectrales du traité réside dans ces
opérateurs et leurs entrelacements. La démonstration de leurs résultats
analytiques ultérieurs conserve son développement propre ; elle n'est pas remplacée par
la périodicité d'une observation résiduelle.

\subsubsection{Discrépance des vacances et lectures de Dirichlet}

Le lien admet également une expression analytique exacte. On conserve le
même codage \(z_n=z_n(0)\), pour \(n\geq1\), et on considère
deux observables ultérieurs : la série sur tous les indices et le
produit sur les indices sélectionnés par irréductibilité. Pour
\(\operatorname{Re}s>1\),
\[
 Z_{\rm vac}^{+}(s)=\sum_{n\geq1}\frac{z_n}{n^s},\qquad
 Z_{\rm vac}^{\times}(s)=\prod_p(1-p^{-s})^{-z_p}.
\]
Tous deux convergent absolument dans ce demi-plan. Leur facteur commun est la
densité \(\delta\), mais leurs discrépances sont différentes.

\Needspace{21\baselineskip}
\begin{proposition}[Séparation de la densité de vacances]
On a
\[
 Z_{\rm vac}^{+}(s)=\delta\zeta(s)+H_{\rm vac}(s),
 \qquad H_{\rm vac}\in\mathcal O(\operatorname{Re}s>0).
\]
Dans \(\operatorname{Re}s>1\), avec les logarithmes définis par les séries absolument convergentes d'Euler,
\[
 Z_{\rm vac}^{\times}(s)=\zeta(s)^\delta G_{\rm vac}(s),\qquad
 \log G_{\rm vac}(s)=
 \sum_p(z_p-\delta)\log(1-p^{-s})^{-1}.
\]
\end{proposition}
\begin{proof}
Les sommes partielles centrées satisfont
\[
 A_N:=\sum_{n=1}^N(z_n-\delta)
 =\lfloor(N+1)\delta\rfloor-N\delta,\qquad |A_N|<1.
\]
La sommation d'Abel représente le terme centré par \(s\int_1^\infty A_{\lfloor x\rfloor}x^{-s-1}\,dx\). L'intégrale converge localement uniformément dans \(\operatorname{Re}s>0\), ce qui prouve l'holomorphie. La seconde identité résulte de la séparation de \(z_p=\delta+(z_p-\delta)\) dans le logarithme du produit, au sein de son domaine de convergence absolue.
\end{proof}

La décomposition rend visible ce qu'ajoute la lecture sur les nombres premiers : elle examine
la même suite de vacances sur un sous-ensemble déterminé par la
factorisation. Le contrôle de la discrépance sur tous les entiers donne
le premier prolongement holomorphe ; le second observable conserve une
discrépance propre sur les nombres premiers. Il s'agit de deux applications de la même
structure discrète, avec des domaines analytiques explicites. Cette articulation
explique leur pertinence dans le noyau arithmétique sans transférer dans cet
article la démonstration des résultats terminaux du traité.
